% =====================================================================
%  Chương 3 — Hệ thống Định giá Bất động sản Việt Nam
% =====================================================================
\chapter{Hệ thống Định giá Bất động sản Việt Nam}
\label{ch:batdongsan}

\section{Phát biểu bài toán}
\label{sec:phatbieu-hou}

\subsection{Hiện tượng thực tế và người dùng}

Thị trường nhà đất Việt Nam thiếu một mức giá tham chiếu công khai. Người mua
dựa vào giá rao của người bán, người bán dựa vào giá rao của hàng xóm, và cả hai
bên đều không có cách kiểm chứng độc lập. Hệ quả là khoảng chênh lệch thương
lượng rất rộng, và bên nào ít thông tin hơn thì chịu thiệt.

Một hệ thống định giá tự động không thay thế được thẩm định giá chính thức ---
việc ấy đòi hỏi khảo sát hiện trạng --- nhưng nó cung cấp được \textbf{mốc tham
chiếu độc lập} dựng từ 30.229 giao dịch thật, cùng với điều mà một con số giá
đơn lẻ không nói được: vay được bao nhiêu, trả góp thế nào, và quanh đó có gì.

\subsection{Phát biểu hình thức}

Cho một hồ sơ bất động sản $\mathbf{x}$ gồm thuộc tính vật lý, vị trí hành chính
và tình trạng pháp lý, hãy học hàm

\begin{equation}
f_{\bm\theta}:\ \mathcal{X}\longrightarrow\mathbb{R}_{+},
\qquad
\hat y = f_{\bm\theta}(\mathbf{x}) \approx \texttt{Price}\ \text{(tỷ VNĐ)},
\label{eq:baitoan-hou}
\end{equation}

\noindent sao cho cực tiểu hoá sai số bình phương trung bình~\eqref{eq:mse} trên
tập huấn luyện. Khác bài toán ở Chương~\ref{ch:tieuduong}, đầu ra ở đây là số
thực liên tục nên không có ngưỡng quyết định nào cần chọn --- nhưng bù lại xuất
hiện một vấn đề mà bài phân lớp không có: \textbf{độ chệch có hệ thống theo vùng
giá trị} (Mục~\ref{sec:chech-hou}).

\begin{table}[H]
\centering
\caption{Đặc tả bài toán định giá bất động sản}
\label{tab:dacta-hou}
\small
\begin{tabularx}{\textwidth}{@{}L{4.2cm}X@{}}
\toprule
\textbf{Hạng mục} & \textbf{Nội dung} \\
\midrule
Hiện tượng thực tế   & Biến động giá bất động sản theo cấu trúc vật lý và vị trí \\
Một quan sát         & Một tin đăng rao bán nhà đất \\
Loại bài toán        & Hồi quy (Regression) \\
Biến mục tiêu        & \texttt{Price}, số thực liên tục, đơn vị tỷ VNĐ \\
Số quan sát          & 30.229 tin đăng \\
Đặc trưng đầu vào    & 14 cột --- 9 định lượng và 5 định danh \\
Số chiều sau mã hoá  & 84 chiều \\
Người dùng cuối      & Người mua, người bán, nhà đầu tư, chuyên viên tín dụng \\
Độ đo quyết định     & MAE (đọc thẳng ra tiền) và $R^2$ (so với đoán trung bình) \\
\bottomrule
\end{tabularx}
\end{table}

\section{Tập dữ liệu và quy trình làm sạch}
\label{sec:dulieu-hou}

\subsection{Nguồn dữ liệu}

Tập \textbf{Vietnam Housing Dataset}~\cite{vnhousing} thu thập từ các trang rao bán bất động sản
Việt Nam, công bố công khai trên Kaggle:

\begin{center}
\small\url{https://www.kaggle.com/datasets/huutri148/vietnam-housing-dataset}
\end{center}

\begin{phathien}{Tệp nguồn có BOM ở đầu}
Tệp CSV bắt đầu bằng dấu Byte Order Mark. Đọc bằng \texttt{encoding="utf-8"}
thông thường khiến tên cột đầu tiên mang thêm ký tự vô hình \texttt{\textbackslash ufeff},
và mọi phép truy cập cột theo tên sau đó đều thất bại --- nhưng thất bại theo
kiểu ``không tìm thấy cột'' chứ không phải ``tệp hỏng'', nên rất dễ đi tìm sai
chỗ. Phải đọc bằng \texttt{encoding="utf-8-sig"}.
\end{phathien}

\subsection{Tỷ lệ dữ liệu khuyết}

\begin{figure}[H]
\centering
\includegraphics[width=0.86\textwidth]{hou_missing_ratio.png}
\caption[Tỷ lệ dữ liệu khuyết theo từng thuộc tính]%
{Tỷ lệ dữ liệu khuyết theo từng thuộc tính của tập bất động sản}
\label{fig:hou-missing}
\end{figure}

\bieudo{Biểu đồ cột ngang, sắp giảm dần theo tỷ lệ khuyết.}

\phanphoi{Trải rất rộng và liên tục --- từ $82{,}6\%$ xuống $11{,}9\%$, không có
khoảng đứt gãy rõ rệt giữa nhóm ``khuyết nhiều'' và nhóm ``khuyết ít''.}

\chitiet{\texttt{Balcony direction} đứng đầu với $82{,}6\%$ --- vượt ngưỡng loại
bỏ $80\%$ nên cột này bị loại khỏi mô hình. \texttt{House direction} khuyết
$70{,}3\%$ nhưng vẫn dưới ngưỡng nên được giữ. Tiếp theo là
\texttt{Furniture state} $46{,}7\%$, \texttt{Access Road} $44{,}0\%$,
\texttt{Frontage} $38{,}3\%$, \texttt{Bathrooms} $23{,}4\%$,
\texttt{Bedrooms} $17{,}1\%$, \texttt{Legal status} $14{,}9\%$ và
\texttt{Floors} $11{,}9\%$. Ba cột \texttt{Address}, \texttt{Area} và
\texttt{Price} đầy đủ hoàn toàn.}

\nhanxet{Việc giữ lại \texttt{House direction} dù khuyết $70{,}3\%$ là một quyết
định có căn cứ: hướng nhà mang ý nghĩa phong thuỷ đáng kể trong định giá tại
Việt Nam, và với biến định danh thì bản thân trạng thái ``không rõ'' cũng là một
hạng mục có ý nghĩa --- tin đăng không khai hướng nhà có thể chính là tin đăng
mà hướng nhà không phải điểm mạnh.}

\subsection{Tách địa chỉ thành hai cấp hành chính}

Cột \texttt{Address} là chuỗi tự do. Bước làm sạch tách nó thành
\texttt{Province} và \texttt{District}. Dữ liệu thật để lại nhiều dạng viết khác
nhau của cùng một địa danh --- ``Hà Nội'', ``HN'', ``Quận Nam Từ Liêm'' xuất hiện
ở cấp tỉnh; ``Hồ Chí Minh giá 2tỷ380'' cho thấy người đăng tin gõ lẫn giá vào ô
địa chỉ. Hệ thống giữ nguyên các biến thể này thay vì hợp nhất thủ công, và xử
lý bằng cơ chế gộp hạng mục hiếm ở bước mã hoá (Mục~\ref{sec:bieudien-hou}).

\section{Phân tích khám phá dữ liệu}
\label{sec:eda-hou}

\subsection{Ba biến người mua thật sự cân nhắc}

\begin{figure}[H]
\centering
\includegraphics[width=\textwidth]{hou_price_area_bedrooms.png}
\caption[Khám phá liên hợp Giá -- Diện tích -- Số phòng ngủ]%
{Khám phá liên hợp Giá -- Diện tích -- Số phòng ngủ: ba phân phối biên ở hàng
trên, ba quan hệ chéo ở hàng dưới}
\label{fig:hou-joint}
\end{figure}

\bieudo{Lưới sáu bảng. Hàng trên là ba phân phối biên --- hai biểu đồ tần suất
kèm đường mật độ và một biểu đồ cột cho biến đếm. Hàng dưới là ba quan hệ ---
biểu đồ phân tán, biểu đồ hộp theo nhóm, và phân phối của một biến dẫn xuất.}

\phanphoi{Ba biến ở ba chế độ phân phối hoàn toàn khác nhau. \texttt{Price} gần
đối xứng, \texttt{Area} lệch phải rất mạnh, \texttt{Bedrooms} là biến đếm có một
đỉnh trội hẳn.}

\chitiet{\texttt{Price} có độ lệch $-0{,}029$ --- gần như đối xứng hoàn hảo ---
với trung vị $5{,}90$ tỷ. \texttt{Area} có độ lệch $3{,}889$, khối lượng dữ liệu
dồn trong khoảng $30$--$100$ m$^2$ với trung vị $56$ m$^2$, đuôi phải kéo tới
$595$ m$^2$. \texttt{Bedrooms} có đỉnh áp đảo tại 3 phòng ($13.703$ tin,
$45{,}3\%$), kế đến 4 phòng ($7.740$) và 2 phòng ($4.290$).

Hàng dưới cho ba con số quan trọng: Giá theo Diện tích có $r=0{,}098$; Giá theo
Số phòng ngủ có $r=0{,}333$ với các hộp chồng lấn rất nhiều; và đơn giá mỗi mét
vuông có độ lệch $2{,}135$ với trung vị $102{,}3$ triệu đồng/m$^2$.}

\begin{phathien}{Ba điều đi ngược trực giác, đọc được từ một hình}
\textbf{(1) Không được biến đổi log cho \texttt{Price} ở tập này.} Phép biến đổi
log là thao tác mặc định khi làm việc với giá nhà, vì giá nhà thị trường lệch
phải mạnh. Nhưng \texttt{Price} ở đây gần đối xứng ($-0{,}029$), nên lấy log sẽ
\emph{tạo ra} độ lệch trái thay vì sửa chữa gì. Giả định phải được đo trước khi
viện dẫn --- xem thêm Mục~\ref{sec:price-cat-bien}.

\textbf{(2) Diện tích gần như không định giá được căn nhà.} Hệ số $0{,}098$ có
nghĩa Diện tích một mình giải thích chưa tới $1\%$ phương sai của giá. Điều này
nghe vô lý cho tới khi nhớ rằng dữ liệu trải khắp cả nước: $56$ m$^2$ ở Cầu Giấy
và $56$ m$^2$ ở Phổ Yên là hai mức giá hoàn toàn khác nhau. Diện tích chỉ có ý
nghĩa \emph{sau khi đã cố định vị trí}.

\textbf{(3) Đơn giá mỗi m$^2$ mới là biến giữ đúng hình dáng mà trực giác mong
đợi} --- lệch phải mạnh với độ lệch $2{,}135$. Đây chính là biến bị \emph{ẩn} đi
khi chỉ nhìn Giá và Diện tích một cách tách rời, và là lời giải thích trực quan
cho việc vì sao \texttt{District} chi phối mô hình tới vậy
(Mục~\ref{sec:importance-hou}).
\end{phathien}

\subsection{Phân phối biến mục tiêu}

\begin{figure}[H]
\centering
\includegraphics[width=0.82\textwidth]{hou_hist_price.png}
\caption[Phân phối tần suất của biến mục tiêu \texttt{Price}]%
{Phân phối tần suất của biến mục tiêu \texttt{Price}}
\label{fig:hou-hist-price}
\end{figure}

\bieudo{Biểu đồ tần suất, 30 khoảng.}

\phanphoi{Gần đối xứng, một đỉnh rộng và bằng phẳng ở giữa, hai đuôi cắt cụt đột
ngột ở cả hai đầu.}

\chitiet{Dữ liệu nằm gọn trong dải $1{,}0$ -- $11{,}5$ tỷ VNĐ, không có một quan
sát nào ngoài dải ấy. Đỉnh trải rộng trong khoảng $4$--$6$ tỷ với chiều cao
$1.700$--$1.870$ tin mỗi khoảng. Hai đầu dừng đột ngột chứ không thoải dần ---
dấu hiệu đặc trưng của dữ liệu đã bị cắt biên.}

\begin{phathien}{\texttt{Price} đã bị cắt biên từ nguồn --- và hệ quả không nhỏ}
\label{sec:price-cat-bien}
Thống kê đầy đủ: nhỏ nhất $1{,}00$, lớn nhất $11{,}50$, trung bình $5{,}8721$,
trung vị $5{,}90$, độ lệch chuẩn $2{,}2119$, độ lệch (skewness) $-0{,}029$, độ
nhọn (kurtosis) $-0{,}8464$. Trung bình xấp xỉ trung vị và độ lệch xấp xỉ $0$ ---
một phân phối gần đối xứng, gần như đều ở phần giữa.

Đây \textbf{không} phải hình dáng của giá bất động sản thật, vốn lệch phải rất
mạnh với một đuôi dài các bất động sản siêu cao cấp. Kết luận: người thu thập dữ
liệu đã lọc bỏ hai đầu trước khi công bố.

Hai hệ quả phải nói rõ:

\begin{enumerate}
  \item \textbf{$R^2$ ở đây khắt khe hơn bình thường.} Mẫu số
        của~\eqref{eq:r2} là phương sai của \texttt{Price}, và phương sai ấy đã
        bị thu hẹp do cắt biên. Cùng một sai số tuyệt đối sẽ cho $R^2$ thấp hơn
        so với trên một tập trải rộng.
  \item \textbf{Hệ thống chỉ áp dụng được trong đúng dải $1$--$11{,}5$ tỷ.} Mô
        hình cây \emph{không ngoại suy} --- đưa vào một biệt thự $50$ tỷ, nó vẫn
        trả về một con số nằm trong dải đã thấy lúc huấn luyện, và con số ấy sai
        hoàn toàn mà không có cảnh báo nào.
\end{enumerate}

Giả định ``giá nhà lệch phải mạnh'' nghe hiển nhiên tới mức suýt trôi qua mà
không ai đo. Nó đúng với thị trường thật và \emph{sai} với tập dữ liệu này.
\end{phathien}

\subsection{Quan hệ giữa diện tích và giá}

\begin{figure}[H]
\centering
\includegraphics[width=0.82\textwidth]{hou_scatter_area_price.png}
\caption[Quan hệ giữa Diện tích và Giá]%
{Biểu đồ phân tán giữa Diện tích và Giá}
\label{fig:hou-scatter}
\end{figure}

\bieudo{Biểu đồ phân tán (scatter plot) với độ trong suốt cao để thấy mật độ
điểm ở vùng chồng chất.}

\phanphoi{Đám mây điểm dày đặc ở góc trái, thưa dần rất nhanh về phía phải; theo
trục dọc thì trải gần như đều khắp dải giá.}

\chitiet{Với diện tích dưới $130$ m$^2$, giá trải gần như kín toàn bộ dải
$1$--$11{,}5$ tỷ --- cùng một mức diện tích ứng với mọi mức giá. Từ $150$ m$^2$
trở lên điểm thưa hẳn nhưng vẫn không hình thành xu hướng đi lên rõ rệt: có căn
$400$ m$^2$ giá $3{,}8$ tỷ nằm cạnh căn $80$ m$^2$ giá $10$ tỷ. Đường giới hạn
trên ở mức $10$--$11{,}5$ tỷ hiện ra thành một vệt ngang --- chính là biên cắt
đã nêu.}

\subsection{Vị trí hành chính và tình trạng pháp lý}

\begin{figure}[H]
\centering
\begin{subfigure}[b]{0.52\textwidth}
  \includegraphics[width=\textwidth]{hou_bar_province.png}
  \caption{Giá trung bình theo Tỉnh/Thành}
  \label{fig:hou-bar-province}
\end{subfigure}
\hfill
\begin{subfigure}[b]{0.45\textwidth}
  \includegraphics[width=\textwidth]{hou_bar_legal.png}
  \caption{Giá trung bình theo Tình trạng pháp lý}
  \label{fig:hou-bar-legal}
\end{subfigure}
\caption[Giá trung bình theo hai biến định danh]%
{Giá trung bình theo Tỉnh/Thành (12 tỉnh dẫn đầu) và theo Tình trạng pháp lý}
\label{fig:hou-bars}
\end{figure}

\bieudo{Hai biểu đồ cột đứng, mỗi cột là giá trung bình của một hạng mục.}

\phanphoi{Hai hình cho hai câu chuyện trái ngược. Bảng bên trái có độ dốc rõ; bảng
bên phải gần như phẳng hoàn toàn.}

\chitiet{Theo tỉnh, dẫn đầu là Hưng Yên $7{,}895$ tỷ (404 tin), rồi Quảng Ninh
$6{,}984$ (110 tin), Bắc Ninh $6{,}402$ (167 tin), Hồ Chí Minh $6{,}214$
(11.778 tin) và Hà Nội $6{,}154$ (10.457 tin). Chênh lệch giữa đầu và cuối bảng
là khoảng $2{,}3$ tỷ.

Theo pháp lý, ba hạng mục gần như bằng nhau: ``Không rõ'' $5{,}9263$,
``Sale contract'' $5{,}8933$, ``Have certificate'' $5{,}8614$ --- chênh lệch chỉ
$0{,}065$ tỷ, tức $1{,}1\%$.}

\begin{phathien}{Trung bình biên che mất cấu trúc có điều kiện}
Hai hình trên gợi ý rằng pháp lý không ảnh hưởng tới giá, và rằng Hưng Yên là
tỉnh đắt nhất Việt Nam. Cả hai kết luận đều sai, và sai vì cùng một lý do:
\emph{trung bình biên trộn lẫn mọi điều kiện khác}.

Hưng Yên đứng đầu không phải vì đất Hưng Yên đắt hơn đất Quận 1, mà vì 404 tin
đăng ở đó tình cờ là các dự án lớn; trong khi trung bình của Hồ Chí Minh bị kéo
xuống bởi hàng nghìn căn nhỏ ở vùng ven. Tương tự, giấy tờ đầy đủ có giá trị
thật, nhưng những căn thiếu giấy tờ trong tập này lại tập trung ở phân khúc cao
hơn, và hai hiệu ứng triệt tiêu nhau khi lấy trung bình.

Mô hình cây \emph{không} bị đánh lừa như vậy vì nó học quan hệ \emph{có điều
kiện}: sau khi đã chia theo quận, nó mới xét tới pháp lý. Bằng chứng là
\texttt{Legal status} vẫn nhận mức quan trọng $0{,}0133$ ở
Mục~\ref{sec:importance-hou} --- nhỏ, nhưng khác $0$ rõ rệt.
\end{phathien}

\subsection{Ma trận tương quan}

\begin{figure}[H]
\centering
\includegraphics[width=0.94\textwidth]{hou_corr_heatmap.png}
\caption[Ma trận tương quan Pearson của các đặc trưng bất động sản]%
{Ma trận tương quan Pearson của các đặc trưng bất động sản}
\label{fig:hou-corr}
\end{figure}

\bieudo{Bản đồ nhiệt đối xứng, thang màu hai chiều từ xanh (âm) qua trắng (không
tương quan) tới đỏ (dương).}

\phanphoi{Khác hẳn ma trận ở Hình~\ref{fig:dia-corr}: ở đây có cả tương quan âm
rõ rệt, và cột \texttt{Price} thì nhạt màu gần như toàn bộ.}

\chitiet{Cặp mạnh nhất là \texttt{Area}--\texttt{Total\_Area} $=0{,}78$ (đúng
theo định nghĩa) và \texttt{Bedrooms}--\texttt{Bathrooms} $=0{,}74$. Có ba tương
quan âm đáng kể: \texttt{Area}--\texttt{Room\_Density} $=-0{,}59$,
\texttt{Area}--\texttt{Floors} $=-0{,}32$ và \texttt{Area}--\texttt{Frontage\_Ratio}
$=-0{,}26$. Với \texttt{Price}, hệ số cao nhất chỉ là $0{,}364$
(\texttt{Bathrooms}), rồi $0{,}333$ (\texttt{Bedrooms}), $0{,}309$
(\texttt{Floors}), $0{,}290$ (\texttt{Total\_Area}) và $0{,}098$
(\texttt{Area}).}

\begin{phathien}{Trần của mô hình tuyến tính đọc được ngay từ ma trận tương quan}
Không đặc trưng định lượng nào tương quan tuyến tính với giá quá $0{,}364$.
Với một mô hình tuyến tính, $R^2$ tối đa đạt được bị chặn bởi chính các hệ số
này --- và Linear Regression thực tế đạt $R^2 = 0{,}3557$
(Bảng~\ref{tab:ketqua-hou}), gần khít với dự đoán.

Đây là một kiểm chứng đẹp: ma trận tương quan tính được \emph{trước khi} huấn
luyện bất kỳ mô hình nào đã báo trước đúng trần hiệu năng của cả một họ thuật
toán. Nó cũng nói luôn rằng muốn vượt trần ấy thì phải đưa vào thông tin mà các
cột định lượng không chứa --- tức là các biến định danh về vị trí, đúng như Thực
nghiệm~3 sẽ đo được ở Mục~\ref{sec:tn3-hou}.

Quan hệ âm \texttt{Area}--\texttt{Floors} $=-0{,}32$ cũng nói một điều thật về
đô thị Việt Nam: nhà mặt đất diện tích nhỏ thì xây cao lên, còn nhà diện tích
lớn thường là nhà vườn hoặc biệt thự thấp tầng.
\end{phathien}

\section{Thiết kế đặc trưng và biểu diễn}
\label{sec:bieudien-hou}

\begin{table}[H]
\centering
\caption{Ba đặc trưng thiết kế cho bài toán bất động sản}
\label{tab:engineered-hou}
\small
\begin{tabularx}{\textwidth}{@{}L{2.9cm}L{4.4cm}X@{}}
\toprule
\textbf{Đặc trưng} & \textbf{Công thức} & \textbf{Căn cứ} \\
\midrule
\texttt{Total\_Area}     & $\texttt{Area}\times\max(\texttt{Floors},1)$ & Tổng diện tích sàn sử dụng, thứ người mua thật sự trả tiền \\
\addlinespace
\texttt{Room\_Density}   & $\dfrac{\texttt{Bedrooms}+\texttt{Bathrooms}}{\max(\texttt{Area},1)}$ & Mật độ phòng --- phân biệt nhà ở gia đình với nhà chia phòng cho thuê \\
\addlinespace
\texttt{Frontage\_Ratio} & $\dfrac{\texttt{Frontage}}{\max(\texttt{Area},1)}$ & Tỷ lệ mặt tiền --- đại lượng quyết định khả năng kinh doanh \\
\bottomrule
\end{tabularx}
\end{table}

Năm biến định danh được mã hoá one-hot với giới hạn 30 hạng mục phổ biến nhất
mỗi biến; các hạng mục hiếm gộp chung vào một nhóm. Không có giới hạn này,
riêng \texttt{District} với hàng trăm giá trị sẽ tạo ma trận thưa khổng lồ và
khiến mô hình học thuộc từng quận hiếm thay vì học quy luật giá. Kết quả:
biểu diễn cuối cùng có \textbf{84 chiều} từ 14 cột ban đầu.

\begin{table}[H]
\centering
\caption{Kích thước ba tập dữ liệu và kết quả mô hình cơ sở}
\label{tab:split-hou}
\small
\begin{tabular}{@{}lrr@{\hspace{1.4cm}}lr@{}}
\toprule
\textbf{Tập} & \textbf{Số tin} & \textbf{Tỷ lệ} & \textbf{Mô hình cơ sở} & \textbf{Giá trị} \\
\midrule
Huấn luyện & 21.764 & 72{,}0\% & MAE  & 1{,}8587 \\
Kiểm tra   &  5.442 & 18{,}0\% & RMSE & 2{,}2265 \\
Ngoại kiểm &  3.023 & 10{,}0\% & $R^2$ & \textbf{$-0{,}0008$} \\
\textbf{Tổng} & \textbf{30.229} & \textbf{100\%} & & \\
\bottomrule
\end{tabular}
\end{table}

\nhanxet{MAE $=1{,}8587$ của mô hình cơ sở đọc thẳng ra tiền: không dùng mô hình
nào thì sai số trung bình khi định giá một căn nhà là gần \textbf{1,86 tỷ VNĐ}.
$R^2 = -0{,}0008$ xấp xỉ $0$ đúng như lý thuyết --- đoán trung vị không giải
thích được chút phương sai nào.}

\section{Thực nghiệm 1 --- Đối đầu năm mô hình}
\label{sec:tn1-hou}

\begin{table}[H]
\centering
\caption{Kết quả năm mô hình hồi quy trên tập kiểm tra, sắp theo $R^2$}
\label{tab:ketqua-hou}
\small
\begin{tabular}{@{}lrrrr@{}}
\toprule
\textbf{Mô hình} & \textbf{MAE} & \textbf{MSE} & \textbf{RMSE} & \textbf{$R^2$} \\
\midrule
Mô hình cơ sở (trung vị)  & 1{,}8587 & 4{,}9573 & 2{,}2265 & $-0{,}0008$ \\
\midrule
XGBoost                   & \textbf{1{,}0493} & \textbf{1{,}9115} & \textbf{1{,}3826} & \textbf{0{,}6141} \\
Random Forest             & 1{,}0818 & 2{,}0184 & 1{,}4207 & 0{,}5925 \\
Support Vector Regression & 1{,}0743 & 2{,}0883 & 1{,}4451 & 0{,}5784 \\
Decision Tree             & 1{,}2412 & 2{,}7519 & 1{,}6589 & 0{,}4444 \\
Linear Regression         & 1{,}2677 & 3{,}1914 & 1{,}7865 & 0{,}3557 \\
\bottomrule
\end{tabular}
\end{table}

\begin{figure}[H]
\centering
\includegraphics[width=\textwidth]{hou_model_comparison.png}
\caption[So sánh năm mô hình hồi quy trên ba độ đo]%
{So sánh năm mô hình hồi quy trên cùng tập kiểm tra}
\label{fig:hou-comparison}
\end{figure}

\nhanxet{Thứ hạng ở đây \textbf{ngược hẳn} với bài toán tiểu đường. Hai mô hình
tổ hợp phức tạp nhất --- XGBoost và Random Forest --- đứng đầu, còn mô hình
tuyến tính đơn giản nhất đứng cuối. Lời giải thích nằm ở quy mô dữ liệu và cấu
trúc quan hệ:

\begin{itemize}
  \item \textbf{21.764 mẫu huấn luyện} so với 552 ở bài toán tiểu đường --- gấp
        gần 40 lần. Mô hình phương sai cao giờ có đủ dữ liệu để phát huy.
  \item \textbf{Quan hệ giá là phi tuyến và có tương tác mạnh}: hiệu ứng của
        diện tích phụ thuộc hoàn toàn vào quận. Cây quyết định biểu diễn được
        tương tác này một cách tự nhiên; mô hình tuyến tính thì không, và trần
        của nó đã được ma trận tương quan báo trước.
\end{itemize}

Bài học phương pháp: \emph{không có thuật toán nào tốt nhất một cách tuyệt
đối}. Cùng một bộ năm thuật toán, chạy trên hai bài toán, cho hai thứ hạng gần
như đảo ngược. Đây là lý do phải chạy đối đầu thay vì chọn theo danh tiếng.}

\subsection{Giá dự đoán so với giá thật}

\begin{figure}[H]
\centering
\includegraphics[width=0.82\textwidth]{hou_pred_vs_true.png}
\caption[Giá dự đoán so với giá thật của mô hình XGBoost]%
{Giá dự đoán so với giá thật của mô hình XGBoost trên tập kiểm tra}
\label{fig:hou-pred-true}
\end{figure}

\bieudo{Biểu đồ phân tán dự báo -- thực tế, kèm đường chéo tham chiếu ứng với dự
báo hoàn hảo.}

\phanphoi{Đám mây điểm bám quanh đường chéo nhưng \textbf{thoải hơn} đường chéo:
ở vùng giá thấp nó nằm phía trên, ở vùng giá cao nó nằm phía dưới.}

\chitiet{Trong khoảng $4$--$7$ tỷ, đám mây bám sát đường chéo và tương đối hẹp.
Từ $8$ tỷ trở lên, gần như toàn bộ điểm rơi xuống dưới đường chéo --- mô hình
định giá hụt. Ở phía dưới $3$ tỷ thì ngược lại, điểm nằm trên đường chéo. Hình
dáng ``chữ S nằm'' này là dấu hiệu kinh điển của hồi quy về trung bình.}

\subsection{Phân tích phần dư}
\label{sec:chech-hou}

\begin{figure}[H]
\centering
\includegraphics[width=\textwidth]{hou_residual_analysis.png}
\caption[Phân tích phần dư của mô hình XGBoost]%
{Phân tích phần dư của mô hình XGBoost --- bốn góc nhìn về cách mô hình sai}
\label{fig:hou-residual}
\end{figure}

\bieudo{Lưới bốn bảng: biểu đồ phân tán phần dư, biểu đồ tần suất phần dư, đồ
thị Q--Q so với phân phối chuẩn, và biểu đồ cột độ chệch theo ngũ phân vị.}

\phanphoi{Ba bảng đầu trông tương đối lành mạnh --- phần dư tập trung quanh $0$,
gần đối xứng. Bảng thứ tư phá vỡ ấn tượng ấy: các cột không dao động ngẫu nhiên
quanh $0$ mà \textbf{giảm đơn điệu} từ dương sang âm.}

\chitiet{Bảng 1: đám mây phần dư nghiêng rõ rệt, dày ở vùng dự báo $4$--$7$ tỷ.
Bảng 2: phân phối phần dư gần đối xứng quanh $0$ với hai đuôi tương đương. Bảng
3: Q--Q bám đường chuẩn ở phần giữa nhưng lệch ở hai đuôi --- phần dư có đuôi
dày hơn phân phối chuẩn. Bảng 4 là bảng quyết định, chi tiết ở
Bảng~\ref{tab:chech-hou}.}

\begin{table}[H]
\centering
\caption{Độ chệch và MAE theo ngũ phân vị giá thật, mô hình XGBoost}
\label{tab:chech-hou}
\small
\begin{tabular}{@{}lrrrrr@{}}
\toprule
\textbf{Ngũ phân vị} & \textbf{Số tin} & \textbf{Giá thật TB} & \textbf{Giá dự báo TB} & \textbf{Độ chệch} & \textbf{MAE} \\
\midrule
Q1 --- rẻ nhất & 1.089 & 2{,}6849 & 3{,}7338 & \textbf{$+1{,}0488$} & 1{,}1575 \\
Q2             & 1.114 & 4{,}5032 & 5{,}1497 & $+0{,}6465$ & 0{,}9366 \\
Q3 --- giữa    & 1.133 & 5{,}9186 & 6{,}1605 & $+0{,}2419$ & \textbf{0{,}7889} \\
Q4             & 1.030 & 7{,}2372 & 6{,}8511 & $-0{,}3860$ & 0{,}8375 \\
Q5 --- đắt nhất& 1.076 & 8{,}9774 & 7{,}4930 & \textbf{$-1{,}4844$} & \textbf{1{,}5335} \\
\bottomrule
\end{tabular}
\end{table}

\begin{phathien}{Độ chệch đơn điệu theo phân khúc --- phát hiện phải nói với người dùng}
Độ chệch không dao động ngẫu nhiên mà giảm đều từ $+1{,}0488$ xuống $-1{,}4844$
qua năm ngũ phân vị. Hai nguyên nhân cộng dồn:

\begin{enumerate}
  \item \textbf{Hồi quy về trung bình~\cite{hastie2009}.} Cực tiểu hoá MSE ở~\eqref{eq:mse} kéo
        mọi dự báo về phía trung tâm phân phối, vì sai lớn bị phạt bình phương.
  \item \textbf{Mô hình cây không ngoại suy.} Giá trị dự báo của một cây là
        trung bình các mẫu huấn luyện trong lá, nên nó không bao giờ vượt quá
        khoảng giá trị đã thấy.
\end{enumerate}

MAE cũng không đều: $0{,}7889$ ở nhóm giữa nhưng $1{,}5335$ ở nhóm đắt nhất ---
gần gấp đôi. Con số MAE gộp $1{,}0493$ là trung bình của hai chế độ sai rất khác
nhau, và nó \emph{che mất} điều đó.

\textbf{Hệ quả bắt buộc phải hiển thị cho người dùng:} với bất động sản cao cấp,
con số hệ thống đưa ra là \textbf{cận dưới}, không phải ước lượng không chệch.
Một hệ thống định giá im lặng về điều này sẽ khiến người bán nhà cao cấp bị
thiệt một cách có hệ thống.
\end{phathien}

\section{Thực nghiệm 2 --- Khảo sát độ sâu của rừng ngẫu nhiên}
\label{sec:tn2-hou}

\begin{figure}[H]
\centering
\includegraphics[width=0.88\textwidth]{hou_depth_sweep.png}
\caption[Ảnh hưởng của độ sâu tối đa tới $R^2$ của Random Forest]%
{Ảnh hưởng của độ sâu tối đa tới $R^2$ của Random Forest}
\label{fig:hou-depth}
\end{figure}

\bieudo{Biểu đồ đường. Trục ngang là \texttt{max\_depth}, trục dọc là $R^2$.}

\phanphoi{Dạng bão hoà cổ điển: tăng nhanh ở đầu dải, chậm dần, rồi đi ngang
hoàn toàn --- không có đoạn đi xuống, tức không thấy dấu hiệu học thuộc.}

\chitiet{$R^2$ tăng từ $0{,}3712$ ở độ sâu $4$ lên $0{,}5654$ ở độ sâu $12$ ---
phần lớn mức cải thiện nằm ở đoạn này. Từ $16$ tới $20$ chỉ còn tăng nhẹ
($0{,}5862 \to 0{,}5937$), rồi đạt đỉnh $0{,}5954$ tại độ sâu $24$. Từ $28$ trở
đi đường cong bằng phẳng tuyệt đối: $0{,}5954$, $0{,}5953$, $0{,}5953$.}

\begin{table}[H]
\centering
\caption{Khảo sát độ sâu tối đa --- kết quả đầy đủ}
\label{tab:depth-hou}
\small
\begin{tabular}{@{}rrrr@{}}
\toprule
\textbf{max\_depth} & \textbf{MAE} & \textbf{RMSE} & \textbf{$R^2$} \\
\midrule
4  & 1{,}4240 & 1{,}7648 & 0{,}3712 \\
8  & 1{,}2275 & 1{,}5506 & 0{,}5146 \\
12 & 1{,}1376 & 1{,}4671 & 0{,}5654 \\
16 & 1{,}0947 & 1{,}4317 & 0{,}5862 \\
20 & 1{,}0774 & 1{,}4186 & 0{,}5937 \\
\textbf{24} & 1{,}0716 & 1{,}4156 & \textbf{0{,}5954} \\
28 & 1{,}0704 & 1{,}4157 & 0{,}5954 \\
40 & 1{,}0703 & 1{,}4159 & 0{,}5953 \\
\bottomrule
\end{tabular}
\end{table}

\nhanxet{Bẫy ``cực trị chạm biên'' xuất hiện lần thứ hai trong dự án này. Dải
khảo sát ban đầu dừng ở $24$ và $R^2$ đạt đỉnh đúng tại đầu mút. Chỉ sau khi nới
dải tới $40$ mới xác nhận được rằng đó là đỉnh thật chứ không phải điểm cắt cụt.

Việc đường cong \emph{không} đi xuống ở độ sâu lớn là nhờ
\texttt{min\_samples\_leaf=5}: tham số này chặn cây mọc tới lá chỉ chứa một quan
sát. Nó còn có tác dụng phụ quan trọng về kỹ thuật --- để mặc định, rừng 200 cây
trên 21.764 dòng chiếm $156$ MB trên đĩa và vượt giới hạn $100$ MB mỗi tệp của
GitHub. Đặt bằng $5$ thì còn $44$ MB mà $R^2$ chỉ giảm từ $0{,}5976$ xuống
$0{,}5925$.}

\section{Thực nghiệm 3 --- Ảnh hưởng của cách biểu diễn dữ liệu}
\label{sec:tn3-hou}

\textbf{Thiết kế.} Ba phương án biểu diễn, cùng chạy trên Random Forest độ sâu
$16$, cùng phép chia dữ liệu.

\begin{table}[H]
\centering
\caption{Ba phương án biểu diễn và kết quả trên Random Forest}
\label{tab:rep-hou}
\small
\begin{tabular}{@{}lrrrr@{}}
\toprule
\textbf{Phương án biểu diễn} & \textbf{MAE} & \textbf{MSE} & \textbf{RMSE} & \textbf{$R^2$} \\
\midrule
Định lượng $+$ định danh, chưa thiết kế & \textbf{1{,}0932} & \textbf{2{,}0294} & \textbf{1{,}4246} & \textbf{0{,}5903} \\
Đầy đủ $+$ đặc trưng thiết kế           & 1{,}0947 & 2{,}0497 & 1{,}4317 & 0{,}5862 \\
Chỉ đặc trưng định lượng                & 1{,}4388 & 3{,}2988 & 1{,}8163 & 0{,}3340 \\
\bottomrule
\end{tabular}
\end{table}

\begin{figure}[H]
\centering
\includegraphics[width=\textwidth]{hou_representation.png}
\caption[Ảnh hưởng của cách biểu diễn dữ liệu tới $R^2$]%
{Ảnh hưởng của cách biểu diễn dữ liệu tới $R^2$ của Random Forest}
\label{fig:hou-representation}
\end{figure}

\bieudo{Biểu đồ cột ngang, sắp giảm dần theo $R^2$.}

\phanphoi{Hai cột đầu gần bằng nhau, cột thứ ba \emph{ngắn hơn hẳn} --- một
khoảng đứt gãy rõ rệt, khác hẳn hình dáng gần đều ở Hình~\ref{fig:dia-representation}.}

\chitiet{Hai phương án có biến định danh cho $R^2$ $0{,}5903$ và $0{,}5862$ ---
chênh nhau $0{,}0041$, tức trong nhiễu. Bỏ hết biến định danh làm $R^2$ tụt
xuống $0{,}3340$ --- mất $0{,}2563$, tức hơn $43\%$ năng lực giải thích. MAE
đồng thời xấu đi từ $1{,}0932$ lên $1{,}4388$ tỷ VNĐ.}

\begin{table}[H]
\centering
\caption{Biên độ do đổi biểu diễn so với biên độ do đổi thuật toán, bài toán bất động sản}
\label{tab:biendo-hou}
\small
\begin{tabular}{@{}lrl@{}}
\toprule
\textbf{Nguồn biến thiên} & \textbf{Biên độ $R^2$} & \textbf{Ghi chú} \\
\midrule
Đổi cách biểu diễn (3 phương án) & 0{,}2563 & $0{,}3340 \to 0{,}5903$ \\
Đổi thuật toán (5 mô hình)       & 0{,}2584 & $0{,}3557 \to 0{,}6141$ \\
\midrule
\textbf{Tỷ số} & \multicolumn{2}{l}{\textbf{Ngang bằng --- nhưng đổi biểu diễn ít công hơn nhiều}} \\
\bottomrule
\end{tabular}
\end{table}

\nhanxet{Hai kết luận, một trong đó đi ngược kỳ vọng ban đầu.

\textbf{(1) Vị trí hành chính là nhóm thông tin quyết định.} Đây là kết luận
mạnh nhất của cả chương, và nó nhất quán với ba bằng chứng độc lập: tương quan
tuyến tính yếu của mọi biến định lượng (Hình~\ref{fig:hou-corr}), mức quan trọng
đặc trưng (Mục~\ref{sec:importance-hou}), và thực nghiệm này.

\textbf{(2) Ba đặc trưng thiết kế \emph{không} giúp gì.} $R^2$ giảm nhẹ từ
$0{,}5903$ xuống $0{,}5862$ khi thêm chúng vào. Lý do giống hệt trường hợp
\texttt{Glucose\_Level} ở Mục~\ref{sec:importance-dia}: cây quyết định
\emph{tự} tìm được các tổ hợp cần thiết, nên cung cấp sẵn
$\texttt{Area}\times\texttt{Floors}$ chỉ thêm một cột tương quan $0{,}78$ với cột
đã có mà không thêm thông tin nào.}

\section{Thực nghiệm 4 --- Kiểm chuẩn ngoài phân phối}
\label{sec:tn4-hou}

\begin{figure}[H]
\centering
\includegraphics[width=0.92\textwidth]{hou_ood_comparison.png}
\caption[So sánh $R^2$ giữa tập kiểm tra và tập ngoại kiểm]%
{So sánh $R^2$ của năm mô hình giữa tập kiểm tra và tập ngoại kiểm}
\label{fig:hou-ood}
\end{figure}

\begin{table}[H]
\centering
\caption{Kết quả trên tập ngoại kiểm (3.023 tin chưa từng dùng trong phát triển)}
\label{tab:ood-hou}
\small
\begin{tabular}{@{}lrrrr@{}}
\toprule
\textbf{Mô hình} & \textbf{MAE} & \textbf{RMSE} & \textbf{$R^2$} & \textbf{$\Delta R^2$ so với kiểm tra} \\
\midrule
XGBoost                   & \textbf{1{,}1070} & \textbf{1{,}4651} & \textbf{0{,}5629} & $-0{,}0512$ \\
Random Forest             & 1{,}1310 & 1{,}4932 & 0{,}5460 & $-0{,}0465$ \\
Support Vector Regression & 1{,}1302 & 1{,}5279 & 0{,}5246 & $-0{,}0538$ \\
Linear Regression         & 1{,}2935 & 1{,}6571 & 0{,}4408 & $+0{,}0851$ \\
Decision Tree             & 1{,}2736 & 1{,}7038 & 0{,}4088 & $-0{,}0356$ \\
\bottomrule
\end{tabular}
\end{table}

\nhanxet{Khác hẳn bài toán tiểu đường, ở đây kết quả ngoại kiểm \textbf{ổn định
và có thể tin được}: thứ hạng bốn mô hình đầu giữ nguyên hoàn toàn, và mức sụt
$R^2$ đều đặn quanh $0{,}05$ cho ba mô hình dẫn đầu. Sự khác biệt đến từ kích
thước tập: $3.023$ tin so với $77$ ca --- gấp gần 40 lần, nên dao động lấy mẫu
nhỏ hơn nhiều.

Mức sụt $0{,}05$ là biên độ lành mạnh: nó xác nhận điểm số trên tập kiểm tra có
lạc quan chút ít (do đã dùng tập ấy để chọn mô hình) nhưng không phải học thuộc.
Riêng Linear Regression \emph{tăng} $R^2$ --- mô hình có độ chệch cao thì ít
nhạy với việc đổi tập, đúng như lý thuyết dự đoán.}

\section{Mức độ quan trọng của đặc trưng}
\label{sec:importance-hou}

\begin{figure}[H]
\centering
\includegraphics[width=0.9\textwidth]{hou_feature_importance.png}
\caption[Mức độ quan trọng của đặc trưng theo XGBoost]%
{Mức độ quan trọng của 15 đặc trưng dẫn đầu theo mô hình XGBoost}
\label{fig:hou-importance}
\end{figure}

\bieudo{Biểu đồ cột ngang, 15 cột one-hot dẫn đầu.}

\phanphoi{Phân bố phân mảnh --- không cột đơn lẻ nào vượt $0{,}07$, nhưng các cột
dẫn đầu gần như đều thuộc cùng hai nhóm gốc.}

\chitiet{Dẫn đầu là \texttt{Province\_Bình Dương} $0{,}069$, rồi
\texttt{Province\_Hồ Chí Minh} $0{,}056$, \texttt{District\_Hoài Đức} $0{,}050$,
\texttt{Province\_Hà Nội} $0{,}046$. Đặc trưng định lượng đầu tiên xuất hiện là
\texttt{Bathrooms} và \texttt{Total\_Area}, cả hai quanh $0{,}028$.}

\begin{table}[H]
\centering
\caption{Mức quan trọng cộng dồn theo nhóm đặc trưng gốc}
\label{tab:importance-grouped}
\small
\begin{tabular}{@{}lrlr@{}}
\toprule
\textbf{Nhóm đặc trưng} & \textbf{Tổng} & \textbf{Nhóm đặc trưng} & \textbf{Tổng} \\
\midrule
\texttt{District}       & \textbf{0{,}4301} & \texttt{Floors}         & 0{,}0122 \\
\texttt{Province}       & \textbf{0{,}4235} & \texttt{Access Road}    & 0{,}0075 \\
\texttt{Bathrooms}      & 0{,}0287 & \texttt{Area}           & 0{,}0072 \\
\texttt{Total\_Area}    & 0{,}0270 & \texttt{Bedrooms}       & 0{,}0047 \\
\texttt{House direction}& 0{,}0217 & \texttt{Frontage\_Ratio}& 0{,}0041 \\
\texttt{Legal status}   & 0{,}0133 & \texttt{Room\_Density}  & 0{,}0038 \\
\texttt{Furniture state}& 0{,}0131 & \texttt{Frontage}       & 0{,}0031 \\
\bottomrule
\end{tabular}
\end{table}

\begin{phathien}{Hai biến vị trí chiếm 85\% toàn bộ năng lực quyết định}
Cộng dồn theo nhóm gốc cho một con số dứt khoát: \texttt{District} $+$
\texttt{Province} $= 0{,}8536$, tức \textbf{85,4\%} tổng mức quan trọng. Mười
hai nhóm đặc trưng còn lại --- gồm toàn bộ thuộc tính vật lý và pháp lý --- chia
nhau $14{,}6\%$ còn lại.

Con số này giải thích trọn vẹn kết quả Thực nghiệm~3: bỏ biến định danh đi thì
mất đúng phần lớn nhất của cấu trúc, và $R^2$ tụt từ $0{,}5903$ xuống $0{,}3340$.

Nó cũng khép lại nghịch lý ở Hình~\ref{fig:hou-bar-legal}: \texttt{Legal status}
có trung bình biên gần như phẳng nhưng vẫn nhận mức quan trọng $0{,}0133$ --- gấp
gần bốn lần \texttt{Frontage}. Mô hình cây khai thác được quan hệ \emph{có điều
kiện} mà trung bình biên không nhìn thấy.

Câu ngạn ngữ ``vị trí, vị trí, vị trí'' của giới bất động sản, ở tập dữ liệu
này, có một con số cụ thể: $85{,}4\%$.
\end{phathien}

\section{Tích hợp Đồ thị Tri thức Đô thị -- Tài chính}
\label{sec:kg-hou}

\subsection{Phân khúc thị trường}

\begin{table}[H]
\centering
\caption{Ba phân khúc thị trường và gói tín dụng tương ứng}
\label{tab:phankhuc-hou}
\small
\begin{tabularx}{\textwidth}{@{}L{3.2cm}L{1.5cm}L{3.4cm}rrX@{}}
\toprule
\textbf{Phân khúc} & \textbf{Ngưỡng} & \textbf{Gói vay} & \textbf{LTV} & \textbf{Lãi suất} & \textbf{Kỳ hạn} \\
\midrule
Cao cấp    & $\ge 8{,}0$ tỷ & Bất động sản Cao cấp & 70\% & 8{,}5\% & 25 năm \\
Trung cấp  & $\ge 4{,}5$ tỷ & Mua nhà Linh hoạt    & 75\% & 7{,}9\% & 20 năm \\
Bình dân   & $\ge 0{,}0$ tỷ & Nhà ở Xã hội         & 80\% & 4{,}8\% & 25 năm \\
\bottomrule
\end{tabularx}
\end{table}

\subsection{Đồ thị của một ca cụ thể}

\begin{figure}[H]
\centering
\includegraphics[width=\textwidth]{kg_property.png}
\caption[Đồ thị Tri thức Đô thị -- Tài chính của tin đăng \#23933]%
{Đồ thị Tri thức Đô thị -- Tài chính của tin đăng \#23933 --- Phân khúc Cao cấp,
20 thực thể và 20 cạnh}
\label{fig:kg-property}
\end{figure}

\bieudo{Đồ thị có hướng vẽ theo bố cục phân tầng, màu nút mã hoá loại thực thể.}

\phanphoi{Cấu trúc hình thoi: mở rộng ở tầng thuộc tính, thắt lại tại nút định
giá và nút phân khúc, rồi mở rộng lần nữa ở tầng tài chính và hạ tầng.}

\chitiet{Bên trái là bốn thuộc tính vật lý cùng nút pháp lý. Ở giữa là chuỗi
Bất động sản $\to$ Quận Nhà Bè $\to$ Tỉnh Hồ Chí Minh và song song là Mô hình
XGBoost $\to$ Giá $8{,}42$ tỷ $\to$ Phân khúc Cao cấp. Bên phải chia hai nhánh:
nhánh tài chính đi qua Gói vay $\to$ Trả góp $47{,}5$ triệu/tháng $\to$ Thu nhập
tối thiểu $118{,}7$ triệu/tháng, và nhánh hạ tầng đi qua Cụm Hạ tầng Đô thị Cao
cấp tới bốn tiện ích.}

\subsection{Gói tài chính mua nhà}
\label{sec:goi-tai-chinh}

Áp dụng công thức niên kim~\eqref{eq:nienkim} và ngưỡng DTI~\eqref{eq:dti} cho
tin đăng \#23933:

\begin{table}[H]
\centering
\caption{Hồ sơ tài chính đầy đủ cho tin đăng \#23933 (giá dự báo 8,42 tỷ VNĐ)}
\label{tab:taichinh-hou}
\small
\begin{tabularx}{\textwidth}{@{}L{5.2cm}R{3.4cm}X@{}}
\toprule
\textbf{Khoản mục} & \textbf{Giá trị} & \textbf{Cách tính} \\
\midrule
Giá định giá            & 8{,}42 tỷ đ    & Mô hình XGBoost \\
Tỷ lệ cho vay (LTV)     & 70\%           & Theo phân khúc Cao cấp \\
Hạn mức vay $P$         & \textbf{5{,}89 tỷ đ} & $8{,}42\times0{,}70$ \\
Vốn tự có tối thiểu     & 2{,}53 tỷ đ    & $8{,}42-5{,}89$ \\
Lãi suất danh nghĩa     & 8{,}5\%/năm    & Mặt bằng công bố \\
Kỳ hạn                  & 25 năm (300 kỳ) & \\
\textbf{Trả góp hằng tháng $M$} & \textbf{47{,}5 triệu đ} & Công thức~\eqref{eq:nienkim} \\
Tổng tiền lãi           & 8{,}34 tỷ đ    & $Mn - P$ \\
\textbf{Thu nhập tối thiểu} & \textbf{118{,}7 triệu đ/tháng} & $M/0{,}40$, công thức~\eqref{eq:dti} \\
\bottomrule
\end{tabularx}
\end{table}

\begin{phathien}{Tổng tiền lãi vượt cả dư nợ gốc}
Với kỳ hạn 25 năm ở lãi suất $8{,}5\%$/năm, tổng tiền lãi là $8{,}34$ tỷ trên
khoản vay gốc $5{,}89$ tỷ --- tức \textbf{142\% dư nợ gốc}. Người vay trả về cho
ngân hàng tổng cộng $14{,}23$ tỷ để mua một căn nhà định giá $8{,}42$ tỷ.

Đây là con số mà bảng quảng cáo gói vay hiếm khi đặt cạnh con số trả góp hằng
tháng, và hệ thống hiển thị nó tường minh. Lịch trả nợ còn cho thấy một điều
nữa: năm đầu tiên trả $498$ triệu tiền lãi nhưng chỉ giảm được $71$ triệu dư
nợ --- phần lãi dồn hết về đầu kỳ, nên trả trước hạn trong vài năm đầu là lúc
tiết kiệm được nhiều nhất.

Con số \textbf{thu nhập tối thiểu $118{,}7$ triệu đồng mỗi tháng} là câu trả lời
thẳng thắn nhất mà hệ thống đưa ra: nó nói cho người dùng biết căn nhà này có
nằm trong tầm với của họ hay không, điều mà con số ``$8{,}42$ tỷ'' đơn thuần
không nói được.
\end{phathien}

\subsection{Tiện ích hạ tầng đô thị}

\begin{table}[H]
\centering
\caption{Cụm hạ tầng tiện ích theo phân khúc, phân loại theo bốn nhóm chức năng}
\label{tab:tienich-hou}
\small
\begin{tabularx}{\textwidth}{@{}L{3.4cm}rXX@{}}
\toprule
\textbf{Cụm} & \textbf{Bán kính} & \textbf{Tiện ích} & \textbf{Nhóm chức năng} \\
\midrule
\multirow{4}{*}{Cao cấp}   & \multirow{4}{*}{2{,}0 km} & Trung tâm thương mại & Thương mại \\
 & & Trường quốc tế        & Giáo dục \\
 & & Bệnh viện đa khoa     & Y tế \\
 & & Tuyến đường vành đai  & Giao thông \\
\midrule
\multirow{4}{*}{Trung cấp} & \multirow{4}{*}{1{,}5 km} & Siêu thị & Thương mại \\
 & & Trường công lập       & Giáo dục \\
 & & Trạm y tế phường      & Y tế \\
 & & Tuyến xe buýt nội đô  & Giao thông \\
\midrule
\multirow{4}{*}{Cơ bản}    & \multirow{4}{*}{3{,}0 km} & Chợ dân sinh & Thương mại \\
 & & Trường tiểu học       & Giáo dục \\
 & & Phòng khám đa khoa    & Y tế \\
 & & Tuyến xe buýt liên huyện & Giao thông \\
\bottomrule
\end{tabularx}
\end{table}

\nhanxet{Bốn nhóm chức năng được giữ cố định qua cả ba phân khúc, chỉ đổi mức
độ. Cách tổ chức này cho phép so sánh trực tiếp: cùng nhu cầu y tế, phân khúc
cao cấp có bệnh viện đa khoa trong bán kính $2$ km còn phân khúc cơ bản có phòng
khám khu vực trong bán kính $3$ km.

Xin nhắc lại ranh giới ở Mục~\ref{sec:ranh-gioi-so-lieu}: đây là mô hình hoá
theo phân khúc, \textbf{không phải} khảo sát thực địa từng quận --- tập dữ liệu
gốc chỉ có địa chỉ hành chính, không có toạ độ, nên không thể suy ra khoảng cách
thật tới một siêu thị cụ thể.}

\section{Ứng dụng web}
\label{sec:app-hou}

\begin{figure}[H]
\centering
\includegraphics[width=0.86\textwidth]{app-desktop-3-dinh-gia.png}
\caption[Giao diện trang Định giá Bất động sản trên máy tính]%
{Giao diện trang Định giá Bất động sản trên máy tính --- biểu mẫu khai thuộc
tính, giá dự báo, hạn mức vay và đồ thị tri thức}
\label{fig:app-hou-desktop}
\end{figure}

\bieudo{Ảnh chụp màn hình bản đã triển khai công khai, khổ máy tính $1440$ px.}

\chitiet{Cột trái là biểu mẫu chia ba khối --- quy mô và kết cấu, vị trí hành
chính, pháp lý và bàn giao --- trong đó danh sách Quận/Huyện được lọc động theo
Tỉnh/Thành đã chọn. Cột phải hiển thị giá dự báo cỡ lớn, thanh phân khúc, đơn
giá mỗi m$^2$ và hạn mức vay. Bốn tab phía dưới: Gói vay, Hạ tầng \& Tài chính,
So sánh, Vì sao.}

\nhanxet{Tab ``Hạ tầng \& Tài chính'' là nơi hệ thống trả lời trọn vẹn câu hỏi
của người mua: bốn thẻ số ở đầu tab cho trả góp hằng tháng, thu nhập tối thiểu,
vốn tự có và tổng tiền lãi; bảng bên trái liệt kê điều khoản vay; bảng bên phải
là lịch trả nợ theo từng năm với cột trả gốc, trả lãi và dư nợ còn lại.

Ứng dụng cũng hiển thị cảnh báo khi tin đăng thiếu giấy tờ: ``chưa có sổ đỏ hoặc
sổ hồng --- phần lớn ngân hàng sẽ không nhận thế chấp tài sản này'', lấy trực
tiếp từ quan hệ \texttt{QUALIFIES\_FOR} trong đồ thị tri thức.}

\begin{figure}[H]
\centering
\includegraphics[width=0.86\textwidth]{app-desktop-5-ha-tang-tai-chinh.png}
\caption[Tab Hạ tầng \& Tài chính --- trả góp, lịch trả nợ và cụm tiện ích]%
{Tab \emph{Hạ tầng \& Tài chính} --- bốn chỉ số tài chính, bảng điều khoản vay,
lịch trả nợ theo năm và cụm tiện ích hạ tầng}
\label{fig:app-hou-tai-chinh}
\end{figure}

\bieudo{Ảnh chụp màn hình tab thứ hai của trang Định giá.}

\chitiet{Bốn thẻ số ở đầu tab: trả góp $47{,}4$ triệu đồng mỗi tháng, thu nhập
tối thiểu $118{,}6$ triệu đồng, vốn tự có $2{,}52$ tỷ và tổng tiền lãi
$8{,}34$ tỷ --- bằng $142\%$ dư nợ gốc. Bảng bên trái liệt kê bảy dòng điều
khoản vay; bảng bên phải là lịch trả nợ theo năm với ba cột trả gốc, trả lãi và
dư nợ cuối năm. Cuối tab là bốn thẻ tiện ích, mỗi thẻ ghi nhóm chức năng và bán
kính tiếp cận.}

\nhanxet{Con số trên giao diện chênh nhẹ so với Bảng~\ref{tab:taichinh-hou}
($47{,}4$ so với $47{,}5$ triệu đồng) vì ứng dụng tính trên giá dự báo đầy đủ
của mô hình, còn bảng trong báo cáo tính trên con số $8{,}42$ tỷ đã làm tròn để
hiển thị. Đây là chênh lệch làm tròn, không phải hai phép tính khác nhau.

Dòng chú thích dưới lịch trả nợ là chỗ hệ thống nói ra điều mà bảng quảng cáo
gói vay thường không nói: năm đầu tiên trả $498$ triệu tiền lãi nhưng chỉ giảm
được $71$ triệu dư nợ.}

\begin{figure}[H]
\centering

\includegraphics[width=0.34\textwidth]{app-mobile-3-dinh-gia.png}
\caption[Giao diện trang Định giá trên điện thoại]%
{Giao diện trang Định giá trên điện thoại, khổ $390$ px}
\label{fig:app-hou-mobile}
\end{figure}
